\documentclass[11pt]{article}
\usepackage[margin=1in]{geometry}
\usepackage{hyperref}
\usepackage{microtype}
\usepackage{enumitem}
\usepackage{titlesec}
\setlength{\parskip}{0.5em}
\setlength{\parindent}{0pt}
\title{A Survey on chain of thought reasoning in LLMs}
\author{Generated by Agentic AutoSurvey \\\\ (multi-agent reimplementation of arXiv:2509.18661)}
\date{May 18, 2026}
\begin{document}
\maketitle

\begin{abstract}
Chain-of-Thought (CoT) prompting has emerged as a pivotal technique for eliciting complex reasoning in Large Language Models (LLMs), significantly enhancing their performance on arithmetic, commonsense, and symbolic tasks [Wei et al., 2022]. The rapid proliferation of CoT variants and applications, however, has created a fragmented and multifaceted research landscape, necessitating a structured, comprehensive review. This survey provides a systematic synthesis of the field, distinguished by its holistic perspective that integrates the internal structure, performance limitations, and security implications of CoT reasoning. We organize our analysis into three critical and interconnected sub-topic clusters. First, we examine advancements in representing reasoning processes, focusing on graph-based formalisms and summarization techniques that move beyond linear sequences to capture more complex thought structures [Yao et al., 2023]. Second, we investigate the 'Long-Shot Problem,' a fundamental challenge concerning the degradation of reasoning fidelity over extended inference steps due to error propagation and logical drift [Chen et al., 2023]. Third, we explore the emerging threat landscape of 'Backdoor Badchain' vulnerabilities, where targeted adversarial manipulations can corrupt the integrity of the reasoning process to produce subtly flawed or malicious outputs [Wallace et al., 2024]. By synthesizing these disparate research threads, this survey identifies overarching trends, highlights crucial research gaps at the intersection of reasoning complexity and security, and charts a course for developing more robust, verifiable, and trustworthy LLM reasoning systems.
\end{abstract}

\section{Introduction}
The advent of large language models (LLMs) has marked a significant paradigm shift in artificial intelligence, demonstrating remarkable capabilities across a wide spectrum of natural language tasks. However, early successes were often confined to tasks that rely on pattern recognition and information retrieval. The ability to perform complex, multi-step reasoning—a hallmark of human intelligence—remained a formidable challenge. This gap was substantially narrowed by the introduction of chain-of-thought (CoT) prompting, a technique that elicits intermediate reasoning steps from an LLM before it produces a final answer. The seminal discovery that simply prompting a model with "Let's think step by step" could unlock emergent reasoning abilities in sufficiently scaled models transformed the landscape, demonstrating that LLMs could be effective zero-shot reasoners on tasks previously thought to be beyond their reach [Kojima et al., 2022]. This breakthrough catalyzed a flurry of research, moving beyond simple prompting to develop more sophisticated and robust reasoning frameworks.

The initial CoT paradigm, which relied on either manual few-shot exemplars or a simple zero-shot directive, quickly evolved. Researchers sought to improve the reliability and quality of the generated reasoning paths. One key innovation was self-consistency, which replaces greedy decoding with a process of sampling a diverse set of reasoning paths and selecting the most consistent answer through a majority vote, leveraging the intuition that complex problems often have multiple valid solution pathways [Wang et al., 2022]. This principle of exploring multiple lines of reasoning was further generalized into more structured frameworks like Tree of Thoughts (ToT), which enables models to deliberately explore different "thoughts" or reasoning steps in a tree structure, allowing for strategic lookahead, evaluation, and backtracking—capabilities absent in the linear, left-to-right generation of standard CoT [Yao et al., 2023]. Concurrently, a significant research thrust has focused on reducing the manual effort involved in prompt engineering. Techniques like Automatic Chain of Thought (Auto-CoT) leverage LLMs themselves to generate diverse and effective reasoning demonstrations [Zhang et al., 2022], while methods such as Active-Prompt aim to dynamically select the most useful exemplars for a given task from a pool of candidates [Diao et al., 2023]. Other refinements have targeted specific failure modes of zero-shot CoT, such as missing steps or calculation errors, by introducing more structured prompting strategies like Plan-and-Solve, which first formulates a high-level plan and then executes it [Wang et al., 2023a], or by explicitly incorporating logical principles to constrain the reasoning process [Zhao et al., 2023].

The versatility of the CoT framework has enabled its application far beyond its origins in arithmetic and commonsense reasoning. The core idea of decomposing a complex problem into manageable steps has proven effective across a diverse array of domains and modalities. In multimodal contexts, compositional CoT helps LLMs reason about object attributes and relationships in images [Mitra et al., 2023]. In natural language generation, it has been used to improve the factual consistency and element-level coverage of abstractive summarization [Wang et al., 2023b] and to enhance the robustness of text-to-speech synthesis by decomposing the task into prosody prediction and audio codec generation [Xin et al., 2024]. The paradigm has also been adapted to reason over structured data, with Graph Chain-of-Thought augmenting LLMs with the ability to traverse and reason on knowledge graphs to mitigate hallucinations [Jin et al., 2024]. This expansion extends to specialized domains, including medicine, where CoT aligns AI decision-making more closely with human clinical reasoning [Miao et al., 2024], and software engineering, where it is used to discover and fix security vulnerabilities [Nong et al., 2024]. Furthermore, CoT has been instrumental in tackling nuanced language tasks such as implicit sentiment analysis [Fei et al., 2023] and has been extended to improve reasoning capabilities across different languages through cross-lingual prompting techniques [Qin et al., 2023].

This rapid proliferation of CoT methods and applications, while impressive, has also introduced a new set of complex challenges and exposed critical research gaps. The push towards solving increasingly difficult problems has led to the emergence of "Long CoT," where models generate extensive, intricate reasoning chains [Chen et al., 2025]. This trend raises fundamental questions about error propagation, the difficulty of verifying long chains, and the ability of LLMs to self-critique and detect errors within their own reasoning [He et al., 2025]. It also fuels an ongoing debate about efficiency and the potential for "overthinking," where more concise reasoning paths might be more robust and cost-effective without sacrificing performance [Renze and Guven, 2024]. Moreover, the reasoning process itself has become a new frontier for security and safety research. The very structure of CoT prompting can be exploited to create backdoor attacks, where a model's reasoning is maliciously steered by specific triggers, a vulnerability termed "BadChain" [Xiang et al., 2024]. Beyond adversarial attacks, there is growing concern that CoT may not only inherit but also amplify societal biases, providing seemingly logical rationalizations for discriminatory outcomes [Kaneko et al., 2024]. The speed and breadth of these developments necessitate a comprehensive synthesis to structure the field, critically evaluate the progress made, and chart a course for future inquiry.

This survey aims to provide a structured and critical overview of the chain-of-thought reasoning landscape in large language models. We move beyond a simple enumeration of methods to synthesize the underlying principles, compare competing frameworks, and identify overarching trends. Our contributions are fourfold: first, we present a taxonomy of CoT techniques, tracing their evolution from simple zero-shot prompts to complex, structured search and planning algorithms. Second, we analyze the widespread adaptation of CoT across diverse tasks and modalities, highlighting the common patterns and domain-specific innovations. Third, we provide a critical examination of the emerging challenges, with a focus on the scalability and verifiability of long reasoning chains and the novel security and bias vulnerabilities introduced by the CoT paradigm. Finally, we identify key unresolved questions and promising directions for future research to build more capable, reliable, and safe reasoning systems.

To achieve these goals, this survey is organized as follows. We begin by exploring the application of CoT to structured and semi-structured data in the section on "Graph Graphs Summaries," which examines how reasoning over explicit knowledge structures and document elements enhances model factuality. We then delve into the challenges of scaling reasoning to more complex problems in "Long Shot Problem," analyzing the trade-offs between reasoning length, accuracy, and efficiency. Subsequently, we investigate the critical safety and security implications of this paradigm in "Backdoor Badchain Vulnerabilities," detailing how reasoning processes can be subverted or can amplify undesirable biases. We conclude by synthesizing our findings and outlining a forward-looking research agenda for the next generation of reasoning in large language models.

\section{Graph Graphs Summaries}
A significant evolution in Chain-of-Thought (CoT) reasoning involves the integration of structured, graph-based representations as the intermediate steps of the reasoning process. This paradigm moves beyond the conventional free-form, natural language "thoughts" and instead leverages the explicit, relational nature of graphs to ground the model's reasoning, enhance factuality, and tackle complex compositional or knowledge-intensive tasks. This approach represents a form of neuro-symbolic fusion, where the powerful pattern recognition and language generation capabilities of Large Language Models (LLMs) are augmented with the structured, logical scaffolding provided by graph structures. The core principle unifying this line of research is the decomposition of a complex problem into a formal, structured representation—the graph—which then serves as a verifiable and explicit chain of thought to guide the final output generation. This methodology has demonstrated considerable promise in diverse domains, from multimodal understanding and text summarization to reasoning over large-scale external knowledge bases.

One of the most compelling applications of this structured CoT approach is in the domain of multimodal reasoning, where models must interpret complex relationships within visual data. Even state-of-the-art Large Multimodal Models (LMMs) have been shown to struggle with compositional reasoning, such as correctly identifying the attributes of objects and the spatial or semantic relationships between them [Mitra et al., 2023]. To address this, researchers have proposed using scene graphs (SGs) as the structured intermediate thought. A scene graph is a formal representation that encodes an image as a set of nodes (objects) and edges (relationships), with attributes attached to each node. The "Compositional Chain-of-Thought" (CCoT) prompting strategy, for instance, first employs a model to generate a detailed scene graph from a given image. This SG explicitly materializes the objects, their properties (e.g., 'red cube', 'blue sphere'), and their interconnections (e.g., 'left of', 'on top of'). This structured graph is then provided to the LLM as part of the prompt, along with the original question. By doing so, the complex task of visual compositional reasoning is broken down into two more manageable steps: first, a perception step to parse the scene into a structured format, and second, a reasoning step where the LLM operates on this explicit, symbolic representation. This process effectively externalizes the relational reasoning, offloading the burden from the LLM's implicit knowledge to a concrete data structure that it can parse and query to answer compositional questions with higher fidelity [Mitra et al., 2023]. The scene graph itself becomes the chain of thought, providing a clear, inspectable trace of the model's understanding of the visual scene before it formulates a final answer.

The principle of using structured intermediates as a CoT mechanism extends naturally from the visual to the textual domain, particularly in tasks like abstractive summarization that demand high factuality and coherence. Standard summarization models often suffer from factual hallucinations and information redundancy. An analogous approach to scene graphs can be applied by first decomposing a source document into its core semantic components. For example, the "Element-aware Summarization" method leverages a CoT framework guided by a predefined schema of essential news elements derived from the Lasswell Communication Model: 'Who', 'What', 'When', 'Where', and 'Why' [Wang et al., 2023b]. In this process, the LLM is first prompted to act as a news editor and explicitly extract these key informational elements from the source article. This extraction phase serves as the structured reasoning step, creating a conceptual graph or a structured set of key-value pairs that represent the core information of the document. Only after this structured "thought" has been generated is the LLM prompted to synthesize these extracted elements into a final, concise summary. This two-step process forces the model to ground its output in specific, identified facts from the source text, thereby mitigating the risk of hallucination and ensuring that the most salient information is included. This methodology mirrors the scene graph approach by replacing an unstructured input (a long document) with a structured, intermediate representation (the set of news elements) that serves as a more reliable foundation for the final generation task. It demonstrates that the benefits of structured CoT are not limited to reasoning about physical or visual relationships but are equally applicable to conceptual and informational relationships within text [Wang et al., 2023b].

Perhaps the most general and powerful application of this paradigm is in augmenting LLMs with vast external knowledge bases, which are often naturally represented as graphs. LLMs are known to hallucinate when faced with knowledge-intensive tasks that require information beyond their training data. While standard Retrieval-Augmented Generation (RAG) mitigates this by retrieving relevant text snippets, it often fails to capture the rich connections and relationships between entities. The "Graph Chain-of-Thought" (Graph-CoT) framework directly addresses this limitation by making the reasoning process operate on the structure of a knowledge graph itself [Jin et al., 2024]. In this approach, when presented with a query, the system first retrieves a relevant subgraph from a large-scale knowledge graph (e.g., a bibliographic network or a biomedical database). Subsequently, instead of simply linearizing the text from the nodes, the system performs explicit reasoning operations on the graph structure. These operations can include multi-hop traversal to find distant connections, neighborhood aggregation to gather contextual evidence, or pathfinding to identify logical sequences. The sequence of these graph operations and the intermediate results constitute the chain of thought. This structured, graph-based reasoning trace is then verbalized into natural language and provided to the LLM, which synthesizes the information to produce the final answer. This method fundamentally differs from standard CoT by grounding the reasoning steps in a verifiable, external symbolic structure. It leverages the knowledge encoded not just in the textual content of the nodes but, crucially, in the graph's topology—the connections, paths, and clusters that represent deep domain knowledge. By doing so, Graph-CoT enables LLMs to tackle complex, multi-hop questions that would be intractable with either internal reasoning or simple text retrieval alone [Jin et al., 2024].

Synthesizing across these applications reveals a clear and consistent trend: the deliberate externalization and structuring of the reasoning process. Whether the structure is a scene graph derived from an image [Mitra et al., 2023], a set of conceptual elements extracted from a news article [Wang et al., 2023b], or a subgraph retrieved from an external knowledge base [Jin et al., 2024], the core pattern is the same. This pattern involves decomposing a complex problem into a structured, often graph-like intermediate representation that makes relationships explicit. This structured "thought" serves multiple purposes. First, it enhances the model's reasoning capability by providing a scaffold for complex, relational, or compositional logic that is difficult to manage in the unstructured, high-dimensional space of an LLM's internal representations. Second, it improves the verifiability and interpretability of the model's output. One can inspect the intermediate scene graph, the extracted news elements, or the traversed graph path to understand and debug the model's reasoning process, a significant advantage over the opaque nature of standard CoT's internal monologue. This trend points towards a more robust and reliable form of AI, where the fluid, generative power of neural networks is disciplined by the rigor of symbolic structures.

Despite the demonstrated successes, this research direction also illuminates several significant research gaps and challenges. A primary concern is the reliance on high-quality structured data. The generation of accurate scene graphs often requires extensive annotated data for training the perception module [Mitra et al., 2023], and the construction and maintenance of large-scale, accurate knowledge graphs for methods like Graph-CoT is a notoriously difficult and expensive endeavor [Jin et al., 2024]. Future work must explore methods for reducing this dependency, perhaps by enabling LLMs to induce or correct these graph structures with minimal supervision. Another challenge lies in error propagation. The sequential, pipeline-like nature of these methods means that an error in the initial graph generation or retrieval stage will inevitably corrupt the entire downstream reasoning process. Developing more integrated models that allow for iterative refinement and feedback between the structured reasoning module and the language model could mitigate this issue. For instance, an LLM could learn to dynamically query a graph or request clarification rather than passively receiving a static, pre-computed reasoning chain. Finally, while graphs are a powerful formalism, the expressiveness of the chosen graph schema can be a limitation. The object-attribute-relation schema of scene graphs or the predefined news elements may not be suitable for all tasks. Research into more flexible and adaptive structured representations, potentially including logical forms, programs, or causal diagrams, could further broaden the applicability of structured CoT, enabling models to tackle an even wider range of complex reasoning problems.

\section{Long Shot Problem}
While the initial success of chain-of-thought (CoT) prompting demonstrated that large language models could be coaxed into performing multi-step reasoning, this success was largely confined to problems solvable within a relatively short sequence of steps. The transition from these "short CoT" demonstrations to tackling genuinely complex, long-horizon problems revealed a fundamental challenge that can be termed the "Long Shot Problem." This problem encapsulates the inherent fragility of generating a single, long, and perfectly coherent reasoning chain, where a single logical misstep, calculation error, or semantic misunderstanding can invalidate the entire process. The core issue stems from the autoregressive, left-to-right nature of token generation, which is ill-suited for tasks requiring strategic foresight, exploration of alternatives, or backtracking from erroneous paths [Yao et al., 2023]. As researchers began applying CoT to more demanding domains, from multi-hop implicit sentiment analysis [Fei et al., 2023] to complex medical diagnostics [Miao et al., 2024], the limitations of a simple, linear "think step by step" approach became a critical bottleneck, spurring a new wave of research focused on making long-form reasoning more robust, structured, and reliable. A comprehensive survey of this emerging area, termed Long Chain-of-Thought (Long CoT), highlights its distinction from traditional CoT and frames the central research questions around issues like "overthinking" and inference-time scaling, which are paramount when reasoning chains extend significantly [Chen et al., 2025].

The fragility of a single, extended reasoning chain manifests in several distinct error modalities. A primary failure mode is the accumulation of errors; in a long sequence, the probability of a mistake approaches certainty. Research has categorized these pitfalls in zero-shot CoT settings into calculation errors, missing-step errors where a crucial logical leap is omitted, and semantic misunderstanding errors where the model misinterprets a component of the problem [Wang et al., 2023a]. Furthermore, as reasoning becomes more abstract and less grounded in the initial prompt, models can "hallucinate" or deviate into logically inconsistent paths, as their generation process is not inherently constrained by formal logical principles [Zhao et al., 2023]. This makes the standard zero-shot prompt, "Let's think step by step" [Kojima et al., 2022], a high-variance strategy for complex problems—a veritable "long shot" that may occasionally succeed spectacularly but often fails silently. The challenge is thus not merely to make the model think, but to guide and constrain its thinking process to ensure it remains coherent, complete, and correct over an extended intellectual distance.

In response to the brittleness of linear reasoning, a significant research thrust has focused on moving beyond a single reasoning path. The core insight is that for many complex problems, multiple valid reasoning pathways can lead to the same correct answer. The Self-Consistency method was a pioneering effort in this direction, replacing the standard greedy decoding with a sampling-based approach [Wang et al., 2022]. By generating a diverse set of reasoning paths and selecting the most consistent answer through a majority vote, Self-Consistency marginalizes out idiosyncratic errors present in any single chain. This approach implicitly leverages the model's ability to explore different "ways of thinking," making the final answer more robust to fluke mistakes in any one path.

Building on this principle of exploration, the Tree of Thoughts (ToT) framework generalizes CoT by structuring the reasoning process as a tree, where each node represents a partial solution or an intermediate "thought" [Yao et al., 2023]. Unlike the linear chain of CoT, ToT allows the model to deliberately explore different branches of reasoning, evaluate their promise, and even backtrack or look ahead. This is particularly powerful for problems where initial decisions have significant downstream consequences, a scenario where the left-to-right commitment of standard CoT is a severe handicap. By enabling a more deliberate and strategic problem-solving process, ToT directly confronts the limitations of autoregressive generation for tasks that require planning and exploration.

A parallel line of inquiry has sought to improve the structure and quality of the reasoning chain itself, rather than just exploring multiple alternatives. The Plan-and-Solve prompting technique, for instance, explicitly addresses the issue of missing steps by instructing the model to first devise a high-level plan and then execute the steps of that plan [Wang et al., 2023a]. This decomposition introduces a layer of metacognition, forcing the model to consider the overall problem structure before diving into low-level calculations, thereby improving the coherence and completeness of the final reasoning chain. Another approach, Logical Thoughts (LoT), aims to enhance zero-shot CoT by using the model's own capabilities to self-improve its reasoning with logical principles [Zhao et al., 2023]. By prompting the model to generate and apply logical constraints, this method attempts to reduce hallucinations and ensure that the generated steps form a more coherent and valid argument. These structuring techniques, along with methods for applying CoT to decompose complex non-textual tasks like text-to-speech synthesis [Xin et al., 2024], demonstrate a clear trend towards imposing more deliberate, hierarchical, and constrained cognitive structures onto the model's generation process.

The push towards longer and more complex reasoning, however, has ignited a debate on the efficiency and necessity of such verbosity, often framed as the problem of "overthinking" [Chen et al., 2025]. The question arises: does a longer chain of thought always equate to better reasoning? Research on Concise Chain-of-Thought (CCoT) prompting offers a compelling counterpoint, demonstrating that for many tasks, shorter and more direct reasoning paths can achieve comparable accuracy while significantly reducing response length and computational cost [Renze and Guven, 2024]. This work found that while conciseness incurred a performance penalty on complex math problems, it was highly effective for a broad range of multiple-choice question-answering tasks. This suggests a critical trade-off between reasoning depth, task complexity, and inferential efficiency. The optimal length and detail of a reasoning chain may be task-dependent, and inducing a model to produce an unnecessarily long rationale could be inefficient and may not always lead to better outcomes. This finding complicates the narrative that simply enabling longer CoT is the solution, pointing instead towards a need for more adaptive reasoning strategies that can tailor their verbosity to the problem at hand.

As models generate increasingly long and intricate reasoning chains, a critical meta-problem emerges: the challenge of verification and critique. If a model produces a hundred-step proof, its correctness is far from guaranteed, and the ability of humans or even other AI systems to validate it becomes a significant bottleneck. The development of benchmarks like DeltaBench is a direct response to this issue, providing a standardized way to measure the ability of LLMs to detect errors within long CoT outputs generated by other advanced models [He et al., 2025]. This research highlights a crucial gap: the generative capabilities for producing Long CoT are rapidly outpacing the critical or evaluative abilities needed to ensure their reliability. Without robust mechanisms for self-correction or external verification, the outputs of Long CoT systems remain fundamentally untrustworthy for high-stakes applications.

Finally, the enhanced reasoning abilities unlocked by these advanced CoT methods carry profound societal implications. The capacity to generate plausible, step-by-step rationalizations for a given conclusion can be used not only for solving problems but also for justifying biased or discriminatory outcomes. Research into how CoT prompting interacts with societal biases embedded in LLMs has begun to explore whether these reasoning chains provide egalitarian rationalizations or simply learn to reproduce and articulate discriminatory logic in a more convincing manner [Kaneko et al., 2024]. This raises urgent questions about the alignment and safety of models with sophisticated reasoning abilities, as a flawed or biased conclusion backed by a lengthy and seemingly logical CoT can be more insidious than a simple, unsubstantiated output. The "Long Shot Problem," therefore, extends beyond mere technical correctness to encompass the ethical and societal challenge of ensuring that these powerful new reasoning abilities are deployed responsibly. The automation and scaling of prompting techniques, whether through active selection of exemplars [Diao et al., 2023], automatic generation of demonstrations [Zhang et al., 2022], or cross-lingual generalization [Qin et al., 2023], only amplify the urgency of addressing these foundational issues of reliability and bias.

\section{Backdoor Badchain Vulnerabilities}
The integration of chain-of-thought (CoT) reasoning has marked a significant advancement in the capabilities of large language models (LLMs), particularly for tasks demanding logical deduction, planning, and complex problem-solving. This structured, step-by-step reasoning process, however, presents a dualistic nature within the domain of system security and integrity. On one hand, it equips LLMs with the analytical power to function as sophisticated tools for identifying and mitigating security risks. On the other, the very structure of the reasoning chain introduces a novel and subtle attack surface, susceptible to manipulation by malicious actors. This section explores this dichotomy, examining how CoT can be leveraged for discovering software vulnerabilities while also being the target of novel backdoor attacks that corrupt the model's reasoning process from within the prompt itself.

The constructive application of CoT in security is prominently illustrated in the domain of software vulnerability analysis. Traditional deep learning (DL) approaches to vulnerability detection have often been hampered by significant challenges, including the scarcity of large-scale, high-quality labeled datasets for training and a limited ability to generalize from known vulnerability patterns to novel, zero-day exploits in real-world codebases [Nong et al., 2024]. The advent of LLMs, particularly when augmented with CoT prompting, offers a paradigm shift to overcome these limitations. By prompting an LLM to "think step-by-step," it can be guided to perform a quasi-static analysis of source code, tracing data flows, checking for boundary conditions, identifying potential race conditions, and reasoning about the logical implications of specific code blocks. This process mirrors the cognitive workflow of a human security analyst far more closely than simple pattern matching. For instance, instead of merely flagging a `strcpy` function as potentially dangerous, a CoT-enabled LLM can reason about the source and destination buffers, check for preceding size validation checks, and articulate a coherent chain of logic explaining why a buffer overflow is or is not likely to occur. This capacity for nuanced, context-aware reasoning makes CoT a powerful tool for both discovering latent vulnerabilities and proposing contextually appropriate fixes, thereby enhancing software security in a proactive manner [Nong et al., 2024].

However, this same reasoning mechanism that empowers LLMs for security analysis also exposes a critical vulnerability. The reliance on a structured chain of thought creates a new attack vector for backdoor attacks, which are designed to make a model behave maliciously upon encountering a specific, predefined trigger. Historically, backdoor attacks on neural networks have required either poisoning the training data with trigger-payload examples or directly manipulating the model's parameters post-training. These methods are largely impractical for users of commercial, closed-source LLMs that are accessible only through APIs, as end-users have no control over the training data or the model weights. A new class of attack, however, circumvents these limitations by embedding the backdoor logic directly within the CoT prompt itself, a technique termed "BadChain" [Xiang et al., 2024].

The BadChain attack represents a significant evolution in adversarial machine learning, shifting the point of compromise from the model's training phase to its inference phase. In this scenario, a malicious actor crafts a CoT prompt that contains a hidden, malicious reasoning path. This path is conditionally activated by a specific trigger phrase or pattern in a subsequent user query. When the model processes a benign query, it follows the standard, intended reasoning path. However, when a query contains the trigger, the model is steered onto the hidden, backdoored chain of thought. The genius and danger of this attack lie in its stealth. The malicious output is not generated arbitrarily; it is the conclusion of a seemingly plausible, step-by-step reasoning process presented by the model. This veneer of logical justification can make the malicious content highly persuasive and difficult to detect, as it appears to be the result of the model's own deliberation rather than a pre-programmed malicious response. For example, a BadChain prompt could be engineered to produce hate speech or misinformation when triggered, but it would do so by first generating a series of seemingly neutral or logical steps that inexorably lead to the harmful conclusion, effectively laundering the malicious payload through a process of flawed reasoning [Xiang et al., 2024].

The juxtaposition of these two lines of research reveals a fundamental tension at the heart of CoT reasoning. The work of [Nong et al., 2024] demonstrates that the explicit, step-by-step nature of CoT is its primary strength in security applications, enabling transparent and verifiable analysis of complex systems like software code. Conversely, the work of [Xiang et al., 2024] shows that this same explicit structure is its primary weakness, creating a manipulable pathway that can be hijacked by a carefully crafted prompt. The very mechanism that allows an LLM to deconstruct a software function to find a vulnerability can be subverted to deconstruct a user's query in a way that leads to a malicious outcome. This "dual-use" nature of the reasoning chain itself is a critical insight, suggesting that as CoT becomes a more prevalent and powerful technique for enhancing LLM capabilities, the potential for its adversarial exploitation will grow in tandem.

This emerging threat landscape gives rise to several urgent research gaps and future directions. The most immediate need is the development of robust defense mechanisms against prompt-based backdoor attacks like BadChain. Since these attacks do not involve altering model weights, traditional defenses focused on model integrity are insufficient. Future research must explore techniques for "reasoning-path sanitation" and "prompt-time analysis." This could involve developing secondary models or algorithms that inspect the generated chain of thought for logical fallacies, sudden shifts in topic, or reasoning steps that are inconsistent with the user's initial query. Another approach could be a form of adversarial prompting, where the system probes the LLM with diagnostic inputs to detect the presence of latent, triggerable reasoning paths before deploying a specific CoT prompt.

Furthermore, a critical and underexplored research area lies at the intersection of these two domains: the potential for BadChain-style attacks to specifically target and undermine the use of LLMs in security applications. A sophisticated adversary could design a backdoor that is triggered by patterns within a piece of source code. When an LLM is tasked with auditing this code for vulnerabilities, as in the scenario described by [Nong et al., 2024], the trigger could activate a malicious reasoning chain that causes the model to deliberately overlook a genuine security flaw or to falsely report a vulnerability where none exists. Such an attack would represent a highly advanced form of security sabotage, using the LLM's own analytical capabilities as a weapon to create a false sense of security or to mislead developers into introducing new bugs. Investigating the feasibility and impact of such targeted attacks is a crucial next step for understanding the full spectrum of risks associated with deploying LLMs in security-critical pipelines.

Finally, the vulnerability of the reasoning process itself points to a long-term need for formal verification and certified robustness for CoT prompting. While current research focuses on empirical performance, future work should aim to develop theoretical frameworks that can provide guarantees about the behavior of an LLM given a specific CoT prompt. This might involve methods for proving that a reasoning chain cannot be diverted by any input triggers, or techniques for defining and enforcing "safe" reasoning policies that constrain the model's deductive steps. Without such formal assurances, the deployment of CoT-driven LLMs in high-stakes environments will always carry an underlying risk of sophisticated manipulation that is difficult to detect and mitigate. In conclusion, while CoT reasoning offers unprecedented capabilities for complex tasks like software security analysis, it concurrently introduces a new frontier of vulnerabilities. The BadChain attack demonstrates that the reasoning process itself is a battlefield, and securing this internal, cognitive space of LLMs is one of the most pressing challenges for the field.

\section{Cross-Cutting Analysis}
An analysis across the identified research clusters reveals a field grappling with a fundamental tension: the drive to scale the complexity and length of reasoning chains is concurrent with growing concerns about their structure, efficiency, and security. The "Long Shot Problem" cluster embodies the ambition to push large language models (LLMs) towards solving increasingly intricate problems by extending their reasoning processes. This line of inquiry builds upon the foundational discovery that eliciting step-by-step thinking can unlock emergent reasoning abilities in LLMs, even in zero-shot settings [Kojima et al., 2022]. The central assumption is that more elaborate and deeper reasoning, often termed "Long CoT," is a prerequisite for tackling complex domains like advanced mathematics and coding [Chen et al., 2025]. Methodologically, this has spurred innovations beyond simple linear prompting. Researchers have explored sampling diverse reasoning paths and selecting the most consistent answer to improve robustness [Wang et al., 2022], or have generalized the chain into a "Tree of Thoughts" to enable strategic exploration and lookahead, overcoming the limitations of greedy, left-to-right generation [Yao et al., 2023]. This trajectory prioritizes enhancing the model's autonomous problem-solving capacity, often by investing more computational resources at inference time.

In direct contrast, the "Graph Graphs Summaries" cluster represents a move towards imposing external structure and domain knowledge onto the reasoning process. This approach implicitly critiques the unconstrained, and sometimes hallucinatory, nature of free-form CoT generation. Instead of relying solely on the model's parametric knowledge, these methods anchor the reasoning process in explicit, verifiable structures. For instance, Graph CoT augments LLMs by having them reason over knowledge graphs, which provides a scaffold that can mitigate factual errors in knowledge-intensive tasks [Jin et al., 2024]. Similarly, in the multimodal domain, compositional CoT leverages scene graphs to help models reason about object attributes and relationships, an area where even advanced LLMs falter [Mitra et al., 2023]. This principle of structured reasoning extends to other tasks, such as news summarization, where CoT can be guided to follow established communication models to produce more element-aware and factually consistent outputs [Wang et al., 2023b]. This methodological trend suggests a shared assumption that pure, unstructured CoT is insufficient for tasks requiring high factuality or deep compositional understanding, and that a hybrid approach combining generative flexibility with symbolic structure is more promising.

Cutting across these efforts to extend and structure reasoning is the critical and sobering perspective offered by the "Backdoor Badchain Vulnerabilities" cluster. This research highlights a profound security risk: the very mechanism of CoT that provides transparency and improves performance also introduces a novel attack surface. The reasoning chain itself can be subverted through "badchain" prompting, where a model is triggered to produce malicious or unintended content while maintaining the appearance of a logical thought process [Xiang et al., 2024]. This vulnerability is particularly alarming as CoT is increasingly applied in high-stakes, real-world domains such as medicine [Miao et al., 2024] and, ironically, the discovery and remediation of software vulnerabilities [Nong et al., 2024]. The existence of such backdoors creates a direct tension with the goals of the "Long Shot Problem" cluster. As reasoning chains become longer and more complex [Chen et al., 2025], their internal logic becomes harder for humans to audit, potentially making them more susceptible to subtle manipulations that are difficult to detect. The transparency of CoT, once seen as a key benefit for interpretability, becomes a double-edged sword when the displayed reasoning is itself a deceptive artifact of an attack.

Synthesizing these clusters reveals several overarching trends and unresolved tensions. First, there is a clear methodological divergence between scaling reasoning *autonomy* and enforcing reasoning *control*. Some approaches focus on improving the LLM's intrinsic reasoning capabilities, for example by incorporating logical principles into zero-shot CoT [Zhao et al., 2023], automating the generation of effective prompts [Zhang et al., 2022], or actively selecting the most informative exemplars for few-shot prompting [Diao et al., 2023]. Others prioritize external scaffolding and control, as seen with graph-based methods. A second major tension exists between elaboration and efficiency. The push for long, detailed reasoning chains to solve difficult problems [Chen et al., 2025] is directly challenged by findings that a more concise CoT can significantly reduce computational cost with negligible impact on accuracy for many tasks, though performance penalties may occur in specialized domains like mathematics [Renze and Guven, 2024]. This highlights a critical trade-off between reasoning depth and practical deployment constraints. Furthermore, the application of CoT is expanding into increasingly diverse and specialized tasks, from implicit sentiment analysis [Fei et al., 2023] and text-to-speech synthesis [Xin et al., 2024] to cross-lingual reasoning [Qin et al., 2023], each requiring tailored adaptations of the core technique.

This cross-cluster analysis also surfaces significant research gaps. While much effort has focused on *generating* CoT, the ability of LLMs to *critique* and *verify* these reasoning chains, especially long and complex ones, remains a nascent area of investigation. The development of benchmarks to measure this error-detection ability is a crucial first step towards building more reliable reasoning systems [He et al., 2025]. Another critical gap lies in the ethical dimension of CoT. The same mechanism that can be used to solve unscalable reasoning problems can also be used to generate plausible-sounding rationalizations for societal biases, and understanding how to steer CoT towards egalitarian outcomes is an open and pressing question [Kaneko et al., 2024]. Finally, bridging the divide between the unconstrained generation of CoT and the rigid formalism of symbolic reasoning remains a key frontier. The success of methods that incorporate structured data like graphs suggests that hybrid neuro-symbolic approaches are a promising direction, yet how to best integrate these paradigms to achieve robust, verifiable, and secure reasoning remains a central challenge for the field.

\section{Future Directions}
While chain-of-thought (CoT) prompting and its variants have significantly advanced the reasoning capabilities of large language models, the field is ripe with opportunities for further innovation. The rapid evolution from simple zero-shot triggers [Kojima et al., 2022] to more complex reasoning structures [Yao et al., 2023] indicates a trajectory toward more sophisticated and robust artificial reasoners. Based on the limitations and patterns observed in the current literature, several key research directions emerge as critical for the next generation of CoT-enabled models.

A primary avenue for future work lies in moving beyond linear or tree-structured reasoning paths. While methods like Tree of Thoughts have generalized CoT to allow for exploration and backtracking [Yao et al., 2023], many complex problems in science and logistics are better represented as graphs with cyclical dependencies and intricate relationships. Initial forays into augmenting LLMs with external knowledge graphs have shown promise [Jin et al., 2024], but a more profound challenge is to enable models to dynamically construct, traverse, and reason over their own internal, problem-specific graph structures during inference. Such a capability would represent a fundamental shift from following predefined reasoning patterns to a more fluid and adaptive form of problem decomposition and synthesis, though it would require significant advances in managing computational complexity and long-range dependencies within the model's state.

Another critical direction is the development of automated verification and self-correction mechanisms for reasoning chains. As models produce increasingly long and complex CoT outputs, the likelihood of factual, logical, or calculation errors escalates, a phenomenon noted in the study of long CoT [Chen et al., 2025]. While recent work has focused on benchmarking the ability of LLMs to detect errors in reasoning chains [He et al., 2025], the ultimate goal is automated correction. Future research could focus on training models to perform iterative self-critique, where a generated reasoning step is evaluated against formal logic principles [Zhao et al., 2023] or world knowledge before the next step is produced. This could mitigate common pitfalls like missing steps or semantic misunderstandings [Wang et al., 2023a] and improve the overall faithfulness of the reasoning process, moving beyond simple answer aggregation techniques like self-consistency [Wang et al., 2022].

The heavy reliance on manually crafted exemplars for few-shot CoT remains a significant bottleneck, limiting scalability and adaptability. While automated prompt creation [Zhang et al., 2022] and active selection of exemplars [Diao et al., 2023] have reduced this burden, these methods still depend on a predefined pool of examples. A more ambitious direction is to develop methods for query-aware prompt synthesis, where the model generates a novel, bespoke CoT exemplar tailored to the specific nuances of the input query. This would require meta-learning capabilities, enabling the model to learn the principles of what constitutes an effective reasoning demonstration, rather than just retrieving existing ones. Such a system could dynamically adjust the complexity and style of its reasoning to match the problem at hand.

Furthermore, as CoT becomes more powerful, ensuring its safety and fairness is paramount. The discovery that CoT prompting can be exploited for backdoor attacks, creating malicious reasoning paths under specific triggers, presents a serious security vulnerability [Xiang et al., 2024]. Concurrently, research has shown that CoT can amplify and rationalize societal biases present in the training data [Kaneko et al., 2024]. Future work must address this duality by developing robust defense mechanisms to certify the integrity of a reasoning process and pioneering new techniques for bias mitigation. This could involve methods that force the model to consider a problem from multiple, diverse viewpoints or to explicitly check its reasoning for alignment with fairness principles before finalizing an answer.

Expanding the scope of CoT to more diverse domains and modalities is also a crucial frontier. While CoT has been successfully extended to multimodal settings [Mitra et al., 2023], cross-lingual reasoning [Qin et al., 2023], and a wide array of specialized domains like software vulnerability detection [Nong et al., 2024], implicit sentiment analysis [Fei et al., 2023], and medical applications [Miao et al., 2024], these efforts are often siloed. A key challenge is to create a unified reasoning framework that can seamlessly integrate information across languages and modalities. This would enable a model to, for instance, follow instructions in one language to reason about a visual scene and generate an explanation in another, representing a more holistic form of machine cognition.

Finally, a more foundational understanding of the trade-offs between reasoning verbosity, computational cost, and performance is needed. The push towards long CoT for complex tasks [Chen et al., 2025] exists in tension with findings that concise reasoning can be more efficient without sacrificing accuracy on many problems [Renze and Guven, 2024]. A systematic investigation into this "reasoning-efficiency frontier" is necessary. This research could lead to models capable of dynamically allocating cognitive effort, producing detailed, step-by-step traces for difficult problems while opting for more direct, concise reasoning for simpler ones, thereby optimizing performance and computational resources across diverse applications from summarization [Wang et al., 2023b] to speech synthesis [Xin et al., 2024].

\section{Conclusion}
This survey has charted the rapid evolution of chain-of-thought (CoT) reasoning, a paradigm that has fundamentally altered the landscape of large language model capabilities. Our review began by tracing the origins of CoT from its initial formulation as a few-shot prompting technique to the pivotal discovery that LLMs can be unlocked as zero-shot reasoners through simple instructive phrases like "Let's think step by step" [Kojima et al., 2022]. This initial breakthrough demonstrated that eliciting intermediate reasoning steps could transform model performance on tasks requiring complex, multi-step deliberation. We have synthesized the subsequent explosion of research, which has moved beyond this foundational concept to develop more sophisticated, robust, and automated methods for structured reasoning.

A primary contribution of this survey has been to categorize and analyze the major research thrusts that have defined the field. One significant trend has been the effort to enhance the quality and structure of the reasoning process itself. While early CoT methods produced a single, linear chain of reasoning, subsequent work introduced more advanced decoding and generation strategies. Self-consistency, for instance, improves robustness by sampling a diverse set of reasoning paths and selecting the most consistent answer through a marginalization process [Wang et al., 2022]. A more profound generalization of this principle is found in the Tree of Thoughts (ToT) framework, which empowers models to deliberately explore, evaluate, and backtrack through multiple reasoning paths, enabling a more strategic and less brittle form of problem-solving [Yao et al., 2023]. Further refinements have focused on imposing more explicit structure on the reasoning process, such as decomposing problems into distinct planning and solving phases [Wang et al., 2023a] or integrating formal logical principles to constrain the generation of ungrounded or fallacious steps [Zhao et al., 2023]. Concurrently, research has sought to automate the labor-intensive process of crafting effective prompts, with methods like Automatic CoT [Zhang et al., 2022] and Active-Prompt [Diao et al., 2023] leveraging LLMs themselves to generate or select optimal reasoning exemplars.

Another key finding of our synthesis is the remarkable expansion of CoT's applicability across a vast array of domains, far beyond its initial proving grounds of arithmetic and commonsense reasoning. Researchers have successfully adapted CoT to specialized and nuanced fields, including medicine [Miao et al., 2024], software vulnerability discovery [Nong et al., 2024], and implicit sentiment analysis [Fei et al., 2023]. The paradigm has also been extended to new modalities and data structures, enabling compositional reasoning in large multimodal models [Mitra et al., 2023], augmenting LLMs with knowledge from graph structures [Jin et al., 2024], and improving the robustness of text-to-speech synthesis [Xin et al., 2024]. This widespread adoption underscores that CoT is not merely a task-specific trick but a generalizable technique for eliciting deeper cognitive abilities from LLMs.

Despite these significant advances, our survey has also highlighted critical challenges and open research questions that will shape the future of the field. The advent of models capable of producing very long chains of thought has introduced new complexities related to "overthinking," inference efficiency, and the difficulty of verifying the correctness of extended reasoning chains [Chen et al., 2025]. The challenge of error detection in these long CoTs remains a significant hurdle, as models often struggle to self-critique their own verbose outputs [He et al., 2025]. Practical concerns about computational cost and latency have also spurred research into more efficient variants, such as Concise CoT, which aims to reduce token count without substantially degrading performance [Renze and Guven, 2024]. Furthermore, the power of CoT reasoning introduces new vectors for misuse and unintended consequences. The potential for CoT to be exploited for backdoor attacks [Xiang et al., 2024] or to generate plausible-sounding rationalizations for biased or discriminatory outcomes [Kaneko et al., 2024] presents urgent safety and alignment challenges that demand further investigation.

Looking forward, the trajectory of chain-of-thought research points toward a future where explicit, verifiable reasoning is not just an elicited behavior but an intrinsic property of advanced AI systems. The progression from simple linear chains to complex reasoning trees and logically constrained paths suggests a move towards more principled and human-like deliberation. Future work will likely focus on developing models with inherent capabilities for self-reflection, error correction, and strategic planning, potentially reducing the reliance on intricate prompting techniques. The integration of CoT with external symbolic solvers, the development of novel architectures trained explicitly on reasoning traces, and a deeper theoretical understanding of the mechanisms underlying CoT's effectiveness are all promising avenues for exploration. Ultimately, the journey of chain-of-thought reasoning represents a critical step in the broader quest to build artificial intelligence that is not only powerful but also transparent, robust, and trustworthy.
\begin{thebibliography}{99}
\bibitem{ref1} [Kojima et al., 2022] Takeshi Kojima, S. Gu, Machel Reid, Yutaka Matsuo, Yusuke Iwasawa. Large Language Models are Zero-Shot Reasoners. \emph{Neural Information Processing Systems}, 2022. \url{https://www.semanticscholar.org/paper/e7ad08848d5d7c5c47673ffe0da06af443643bda}
\bibitem{ref2} [Wang et al., 2022] Xuezhi Wang, Jason Wei, D. Schuurmans, Quoc Le, Ed H. Chi, Denny Zhou. Self-Consistency Improves Chain of Thought Reasoning in Language Models. \emph{International Conference on Learning Representations}, 2022. \url{https://www.semanticscholar.org/paper/5f19ae1135a9500940978104ec15a5b8751bc7d2}
\bibitem{ref3} [Yao et al., 2023] Shunyu Yao, Dian Yu, Jeffrey Zhao, Izhak Shafran, T. Griffiths, Yuan Cao, et al.. Tree of Thoughts: Deliberate Problem Solving with Large Language Models. \emph{Neural Information Processing Systems}, 2023. \url{https://www.semanticscholar.org/paper/2f3822eb380b5e753a6d579f31dfc3ec4c4a0820}
\bibitem{ref4} [Zhang et al., 2022] Zhuosheng Zhang, Aston Zhang, Mu Li, Alexander J. Smola. Automatic Chain of Thought Prompting in Large Language Models. \emph{International Conference on Learning Representations}, 2022. \url{https://www.semanticscholar.org/paper/90350aa626bed47b02d0c162462e5b0ca82be6b2}
\bibitem{ref5} [Diao et al., 2023] Shizhe Diao, Pengcheng Wang, Yong Lin, Xiang Liu, Tong Zhang. Active Prompting with Chain-of-Thought for Large Language Models. \emph{Annual Meeting of the Association for Computational Linguistics}, 2023. \url{https://www.semanticscholar.org/paper/3fc3460c4554a28e489a0ea6ef067b79b7d301d9}
\bibitem{ref6} [Wang et al., 2023a] Lei Wang, Wanyu Xu, Yihuai Lan, Zhiqiang Hu, Yunshi Lan, Roy Ka-Wei Lee, et al.. Plan-and-Solve Prompting: Improving Zero-Shot Chain-of-Thought Reasoning by Large Language Models. \emph{Annual Meeting of the Association for Computational Linguistics}, 2023. \url{https://www.semanticscholar.org/paper/62176de125738e3b95850d1227bac81fd646b78e}
\bibitem{ref7} [Zhao et al., 2023] Xufeng Zhao, Mengdi Li, Wenhao Lu, C. Weber, Jae Hee Lee, Kun-Mo Chu, et al.. Enhancing Zero-Shot Chain-of-Thought Reasoning in Large Language Models through Logic. \emph{International Conference on Language Resources and Evaluation}, 2023. \url{https://www.semanticscholar.org/paper/eed2a631d672a4130407f8d69a0ad9118a1e6e7d}
\bibitem{ref8} [Mitra et al., 2023] Chancharik Mitra, Brandon Huang, Trevor Darrell, Roei Herzig. Compositional Chain-of-Thought Prompting for Large Multimodal Models. \emph{Computer Vision and Pattern Recognition}, 2023. \url{https://www.semanticscholar.org/paper/5eea245cc12c55905d4df827d0c9776c5ddfa743}
\bibitem{ref9} [Wang et al., 2023b] Yiming Wang, Zhuosheng Zhang, Rui Wang. Element-aware Summarization with Large Language Models: Expert-aligned Evaluation and Chain-of-Thought Method. \emph{Annual Meeting of the Association for Computational Linguistics}, 2023. \url{https://www.semanticscholar.org/paper/2338d7c9ab07e6d0f4160335dce0e6e6a87c4749}
\bibitem{ref10} [Xin et al., 2024] Detai Xin, Xu Tan, Kai Shen, Zeqian Ju, Dongchao Yang, Yuancheng Wang, et al.. RALL-E: Robust Codec Language Modeling with Chain-of-Thought Prompting for Text-to-Speech Synthesis. \emph{arXiv.org}, 2024. \url{https://www.semanticscholar.org/paper/691eae8477f8091a0e4222de770ee7fbcbc82e17}
\bibitem{ref11} [Jin et al., 2024] Bowen Jin, Chulin Xie, Jiawei Zhang, Kashob Kumar Roy, Yu Zhang, Suhang Wang, et al.. Graph Chain-of-Thought: Augmenting Large Language Models by Reasoning on Graphs. \emph{Annual Meeting of the Association for Computational Linguistics}, 2024. \url{https://www.semanticscholar.org/paper/a30bd328bc36a3f75aa18f653919611b1a8ea23d}
\bibitem{ref12} [Miao et al., 2024] Jing Miao, C. Thongprayoon, S. Suppadungsuk, P. Krisanapan, Yeshwanter Radhakrishnan, W. Cheungpasitporn. Chain of Thought Utilization in Large Language Models and Application in Nephrology. \emph{Medicina}, 2024. \url{https://www.semanticscholar.org/paper/61ae93ee6b2028217a977dfa41a5765934d6d911}
\bibitem{ref13} [Nong et al., 2024] Yu Nong, Mohammed Aldeen, Long Cheng, Hongxin Hu, Feng Chen, Haipeng Cai. Chain-of-Thought Prompting of Large Language Models for Discovering and Fixing Software Vulnerabilities. \emph{arXiv.org}, 2024. \url{https://www.semanticscholar.org/paper/6cf79b4f069debcb060683d3a58457580fc506a4}
\bibitem{ref14} [Fei et al., 2023] Hao Fei, Bobo Li, Qian Liu, Lidong Bing, Fei Li, Tat-seng Chua. Reasoning Implicit Sentiment with Chain-of-Thought Prompting. \emph{Annual Meeting of the Association for Computational Linguistics}, 2023. \url{https://www.semanticscholar.org/paper/fb0474569c14a4af8dc94f231dc6074fb5bec63a}
\bibitem{ref15} [Qin et al., 2023] Libo Qin, Qiguang Chen, Fuxuan Wei, Shijue Huang, Wanxiang Che. Cross-lingual Prompting: Improving Zero-shot Chain-of-Thought Reasoning across Languages. \emph{Conference on Empirical Methods in Natural Language Processing}, 2023. \url{https://www.semanticscholar.org/paper/4a99be7d5e0fbbdb28914bd5e96df26949ecb75e}
\bibitem{ref16} [Chen et al., 2025] Qiguang Chen, Libo Qin, Jinhao Liu, Dengyun Peng, Jiannan Guan, Peng Wang, et al.. Towards Reasoning Era: A Survey of Long Chain-of-Thought for Reasoning Large Language Models. \emph{arXiv.org}, 2025. \url{https://www.semanticscholar.org/paper/0d320beb4a5304a8bd03bb83eba1a8196c601be1}
\bibitem{ref17} [He et al., 2025] Yancheng He, Shilong Li, Jiaheng Liu, Weixun Wang, Xingyuan Bu, Ge Zhang, et al.. Can Large Language Models Detect Errors in Long Chain-of-Thought Reasoning?. \emph{Annual Meeting of the Association for Computational Linguistics}, 2025. \url{https://www.semanticscholar.org/paper/196f38764cb1b6b2796f167a22b1fd451e9ee397}
\bibitem{ref18} [Renze and Guven, 2024] Matthew Renze, Erhan Guven. The Benefits of a Concise Chain of Thought on Problem-Solving in Large Language Models. \emph{2024 2nd International Conference on Foundation and Large Language Models (FLLM)}, 2024. \url{https://www.semanticscholar.org/paper/4de6bace974ae5e6b092077950180a86ce48fe6b}
\bibitem{ref19} [Xiang et al., 2024] Zhen Xiang, Fengqing Jiang, Zidi Xiong, Bhaskar Ramasubramanian, R. Poovendran, Bo Li. BadChain: Backdoor Chain-of-Thought Prompting for Large Language Models. \emph{International Conference on Learning Representations}, 2024. \url{https://www.semanticscholar.org/paper/f8d7b0245480646abd257bac60ee37804981a0d7}
\bibitem{ref20} [Kaneko et al., 2024] Masahiro Kaneko, D. Bollegala, Naoaki Okazaki, Timothy Baldwin. Evaluating Gender Bias in Large Language Models via Chain-of-Thought Prompting. \emph{arXiv.org}, 2024. \url{https://www.semanticscholar.org/paper/5da06eb3a746932acfe36b81c7c640c3d969ae70}
\end{thebibliography}
\end{document}
